% 11_amorphous.tex -- Part 11, amorphous films: the registered protocol. Sources: research_monograph/chapters/
% M11_amorphous_protocol.tex (sec:mono-am-why, -scope, -mace, -slabs, tab:mono-am-cells, sec:mono-am-mlval,
% tab:mono-am-mace, -relax, -smearing, -restart, -final, -defs, sec:mono-am-hn, tab:mono-am-hn-controls,
% prop:mono-hn, sec:mono-am-a37, sec:mono-am-release); M06_defect_theory.tex (sec:mono-defect-pr,
% tab:mono-defect-pr); M03_literature_materials.tex (sec:mono-lit-amorph); appendices/MG_corrections.tex (X8, X13);
% evidence/CLAIM_SOURCE_MAP.tsv (M001, M005, C034-C037); evidence/FACT_REGISTRY.tsv (FACT-014, FACT-070..073,
% FACT-081); evidence/derived/defect_estimates.json (PR_support).
% Blind rules: protocol facts and force-only checks only; no electronic quantity of any aIO_* / aIWO_* cell;
% per-cell outcomes are \sealed{}. No figure of amorphous electronic structure exists or is used.
% Note: frame bodies start with a command, never with a brace group (beamer would read a subtitle).
\providecommand{\elevenfext}{f_{\mathrm{ext}}}
\providecommand{\elevenflocW}{f_{\mathrm{locW}}}
\providecommand{\elevenDW}{\Delta_{\mathrm{W}}}
\providecommand{\elevenNex}{N_{\mathrm{ex}}}
\providecommand{\elevenNeff}{N_{\mathrm{eff}}}
\providecommand{\elevenEh}{E_{\mathrm{host}}}
\providecommand{\elevenAng}{\mbox{\AA}}
\providecommand{\elevenRyb}{\,Ry/bohr}
% denser slides: footnotesize bullets, local to the frame in which it is called
\providecommand{\elevendense}{\setbeamerfont{itemize/enumerate body}{size=\footnotesize}%
  \setbeamerfont{itemize/enumerate subbody}{size=\footnotesize}}

\partslide{Part 11}{Amorphous films: the registered protocol}

% ---------------------------------------------------------------------------------------------------------------------
\begin{frame}{The channels are disordered; all DFT so far is crystalline}
  \elevendense
  \begin{columns}[T,onlytextwidth]
    \begin{column}{0.56\textwidth}
      \begin{itemize}
        \item The channels are sputtered, largely amorphous or nanocrystalline: HRTEM shows no long-range
          crystallinity\sref[p.~2]{Januar2026}.
        \item Amorphous \InO{} keeps edge- and corner-sharing InO$_x$ polyhedra whose connectivity is linked to
          transport\sref{Buchholz2014}; its s-like conduction band survives disorder\sref{Medvedeva2017}.
        \item Test H-N: does W still put its electron into an extended host state at 0.7\,nm? Confinement, weaker
          screening, surfaces (every atom within about 0.35\,nm of one) and 2-D disorder could prevent it
          \lab{theory}.
        \item Below a mobility edge states are localised; in 2-D weak disorder localises all states on long
          scales\srefs{Anderson1958,Mott1969,Abrahams1979}. A 1--3\,nm cell sees only few-atom localisation, the
          scale of a deep, d-like W trap.
      \end{itemize}
    \end{column}
    \begin{column}{0.40\textwidth}
      \centering
      \begin{tikzpicture}[x=1mm,y=1mm,font=\scriptsize,line width=0.5pt]
        % (top) periodic supercell: three images of one cell of lateral size L
        \foreach \c in {0,16,32,48} {\draw[rulegrey,dashed] (\c,0) -- (\c,8);}
        \draw (0,0) -- (48,0);
        \draw[line width=0.7pt] plot[domain=0:48,samples=220]
          (\x,{7*exp(-((\x-8)^2)/1.4)+7*exp(-((\x-24)^2)/1.4)+7*exp(-((\x-40)^2)/1.4)});
        \draw[line width=0.7pt,inkgrey] (0,1.8) -- (48,1.8);
        \node[anchor=west,inner sep=0.6mm] at (9.4,6.0) {localised};
        \node[anchor=west,inner sep=0.6mm,text=inkgrey] at (25.4,3.1) {extended};
        \draw[{Stealth[length=1.2mm]}-{Stealth[length=1.2mm]}] (16,-1.4) -- node[below,inner sep=0.5mm]{cell, $L = 1$--3\,nm} (32,-1.4);
        \node[anchor=south west,inner sep=0] at (0,8.6) {supercell, periodic images};
        % (bottom) the same disorder in an infinite film: a state localised on a length xi >> L
        \begin{scope}[yshift=-19mm]
          \draw (0,0) -- (48,0);
          \draw[line width=0.7pt] plot[domain=0:48,samples=120] (\x,{7*exp(-abs(\x-20)/10)});
          \draw[ntumaroon,line width=0.7pt] (17.8,7.9) -- (17.8,8.9) -- (22.2,8.9) -- (22.2,7.9);
          \node[anchor=south,text=ntumaroon,inner sep=0.5mm] at (20,8.9) {$L$};
          \draw[{Stealth[length=1.2mm]}-{Stealth[length=1.2mm]}] (10,-1.4) -- node[below,inner sep=0.5mm]{$\xi$ = tens of nm} (30,-1.4);
          \node[anchor=south west,align=left,inner sep=0] at (31.0,3.4) {flat within\\one cell: read\\as extended};
          \node[anchor=south west,inner sep=0] at (0,8.6) {film};
        \end{scope}
      \end{tikzpicture}\par
      \raggedright
      \src{Schematic of $|\psi|^2$ along the film, not to scale. Top: a 1--3\,nm periodic cell separates a
        few-atom state from an extended one. Bottom: a state localised on tens of nm is nearly flat across one
        cell, so the cell reads it as extended.}
    \end{column}
  \end{columns}
\end{frame}

% ---------------------------------------------------------------------------------------------------------------------
\begin{frame}{Every cell goes from ML melt--quench to a certified read-out}
  \elevendense
  \begingroup\centering
  \begin{tikzpicture}[x=1mm,y=1mm,font=\scriptsize,
      st/.style={draw,line width=0.45pt,align=center,text width=17.6mm,minimum height=8.4mm,inner sep=0.8mm},
      ar/.style={-{Stealth[length=1.6mm]},line width=0.5pt}]
    \node[st] (a) at (11,0)   {MACE-MPA-0\\melt--quench};
    \node[st] (b) at (34.6,0) {force check\\against PBE};
    \node[st] (c) at (58.2,0) {PBE relaxation\\(GPU)};
    \node[st] (d) at (81.8,0) {final SCF and\\projections (CPU)};
    \node[st] (e) at (105.4,0) {frozen read-out\\and certificate};
    \node[st,dashed] (f) at (129,0) {release\\(D-43)};
    \foreach \u/\v in {a/b,b/c,c/d,d/e,e/f} {\draw[ar] (\u) -- (\v);}
    \node[anchor=north,text=ntumaroon,inner sep=0.8mm] at (b.south) {failed at W};
    \node[anchor=north,text=inkgrey,inner sep=0.8mm] at (f.south) {sealed};
  \end{tikzpicture}\par\endgroup
  \vspace{0.6mm}
  \begin{columns}[T,onlytextwidth]
    \begin{column}{0.49\textwidth}
      {\footnotesize\textbf{Melt--quench}, MACE-MPA-0 ``medium''\srefs{Batatia2022,Batatia2023}\par}
      \begin{itemize}
        \item Langevin, 1\,fs, friction 0.01\,fs$^{-1}$: 3000\,K for 6\,ps, quench to 300\,K over 9\,ps, hold
          2\,ps; open the vacuum gap, 800\,K for 3\,ps, cool over 3\,ps; relax to 0.05\,eV/\elevenAng, passivate,
          relax again.
        \item Gate on density, bond lengths and coordination; a failed structure is replaced by the next seed
          (three were), thresholds unchanged.
        \item 23\,ps is about $10^{14}$\,K/s: a network of the right composition and density, not a model of the
          sputtered film.
      \end{itemize}
    \end{column}
    \begin{column}{0.48\textwidth}
      {\footnotesize\textbf{PBE relaxation and final SCF}\par}
      \begin{itemize}
        \item 60/480\,Ry; BFGS to $10^{-3}$\elevenRyb{} and $10^{-4}$\,Ry; $2{\times}2{\times}1$ $k$ grid; dipole
          correction. Converged only on ``bfgs converged'' (A19).
        \item Final SCF at 71/568\,Ry with $1.3\times$ the occupied bands, then L\"owdin projections.
        \item The frozen read-out (D-32; 221 tests) must first reproduce two crystal values: a 3\,\% W-5d share and
          $\dEc = 0.281$\,eV at 2\,nm.
      \end{itemize}
    \end{column}
  \end{columns}
\end{frame}

% ---------------------------------------------------------------------------------------------------------------------
\begin{frame}{Six W cells span bulk to 0.7\,nm at one W volume density}
  \elevendense
  \begingroup\scriptsize
  \setlength{\tabcolsep}{5pt}
  \renewcommand{\arraystretch}{1.06}
  \begin{tabularx}{\textwidth}{@{}l l l r l >{\raggedright\arraybackslash}X@{}}
    \toprule
    $t_{\mathrm{eq}}$ & W-doped cell & composition & $N$ & undoped partner & role in H-N \\
    \midrule
    bulk   & \cell{aIWO_bulk_s1} & In$_{49}$W$_1$O$_{76}$, no surface & 126 & \cell{aIO_bulk_s1} & bulk point \\
    2.0\,nm & \cell{aIWO_2p0_s1} & In$_{98}$W$_2$ film, OH/H & 270 & \cell{aIO_2p0_s1} (268) & thick endpoint, seed 1 \\
    2.0\,nm & \cell{aIWO_2p0_s2} & In$_{98}$W$_2$O$_{156}$H$_8$ & 264 & \cell{aIO_2p0_s2} & thick endpoint, seed 2 \\
    1.0\,nm & \cell{aIWO_1p0_s1} & In$_{49}$W$_1$O$_{80}$H$_8$ & 138 & \cell{aIO_1p0_s1} (134) & intermediate \\
    0.8\,nm & \cell{aIWO_0p8_s2} & In$_{49}$W$_1$O$_{76}$ film, OH/H & 147 & \cell{aIO_0p8_s1} & intermediate \\
    0.7\,nm & \cell{aIWO_0p7_s1} & In$_{49}$W$_1$O$_{83}$H$_{14}$ & 147 & \cell{aIO_0p7_s1} & thin decisive endpoint \\
    \bottomrule
  \end{tabularx}\par\endgroup
  \src{$N$: atoms (register and claims ledger); one melt--quench sample per name. 1.5\,nm cut under the budget rule.
    Controls: a compensated 2.0\,nm cell (In$_{98}$W$_2$O$_{160}$H$_{14}$, 274 atoms, no excess electron), bare
    surfaces at 2.0 and 0.7\,nm, \AlO{} and \HfO{} interfaces. Slab rules registered before any ML or DFT result
    (D-14). Structural facts only.}
  \vspace{1.2mm}
  \begin{columns}[T,onlytextwidth]
    \begin{column}{0.49\textwidth}
      \begin{itemize}
        \item W/(In+W) $= 2.0$\,\%. Each W brings one extra O, so one excess electron per W (W$^{6+}$ for In$^{3+}$
          gives three, the O$^{2-}$ takes two): about $6\times10^{20}$\cmcu.
        \item In-plane edge $\ge 12.5$\,\elevenAng{} at 7.0\,g/cm$^3$: the same W density at every thickness.
      \end{itemize}
    \end{column}
    \begin{column}{0.48\textwidth}
      \begin{itemize}
        \item $t_{\mathrm{eq}}$ = film mass / (7.0\,g/cm$^3$ $\times$ area), not the H--H thickness of Parts 6--8.
        \item OH and H in pairs (electron count unchanged); 18\,\elevenAng{} vacuum with a dipole
          correction\sref{Bengtsson1999}.
      \end{itemize}
    \end{column}
  \end{columns}
\end{frame}

% ---------------------------------------------------------------------------------------------------------------------
\begin{frame}{The ML potential fails at W; its structures only seed DFT}
  \elevendense
  \begin{columns}[T,onlytextwidth]
    \begin{column}{0.59\textwidth}
      \begingroup\scriptsize
      \setlength{\tabcolsep}{2.3pt}
      \renewcommand{\arraystretch}{1.08}
      \begin{tabular}{@{}l l c c c@{}}
        \toprule
        structure & frame & RMSE & MAE & RMSE$_{\mathrm{W1}}$ \\
        \multicolumn{2}{@{}l}{\color{inkgrey}registered limit (eV/\elevenAng)} & {\color{inkgrey}0.15} &
          {\color{inkgrey}0.10} & {\color{inkgrey}0.25} \\
        \midrule
        a-\InO, bulk & quenched / final & 0.068 / 0.065 & 0.052 / 0.049 & -- \\
        a-IWO, bulk & quenched / final & 0.121 / 0.122 & 0.068 / 0.068 & \textbf{0.434 / 0.449} \\
        a-IWO, 0.7\,nm & final & \textbf{0.206} & \textbf{0.110} & \textbf{0.815} \\
        melt frame & not scored & 0.311 & 0.171 & 1.108 \\
        \bottomrule
      \end{tabular}\par\endgroup
      \src{MACE forces against PBE single points (QE 7.5, MV 0.01\,Ry) on the same geometries; W1: W and its
        first-shell O; bold: limit exceeded. Comparator frozen before any output was read. C037 \lab{computed}.}
    \end{column}
    \begin{column}{0.38\textwidth}
      \begin{itemize}
        \item Verdict: \textbf{failed overall}. 71--75\,\% of the squared error sits in the W first shell (7 atoms);
          outside it the RMSE is 0.063--0.066 (bulk), 0.107\,eV/\elevenAng{} (0.7\,nm).
        \item Registered consequence: DFT starting points, not validated amorphous models; no ML-stage structural
          statistic is reported (X8).
      \end{itemize}
    \end{column}
  \end{columns}
  \vspace{1.6mm}
  \begingroup\scriptsize
  \setlength{\tabcolsep}{4pt}
  \renewcommand{\arraystretch}{1.08}
  \begin{tabularx}{\textwidth}{@{}l l l l >{\raggedright\arraybackslash}X@{}}
    \toprule
    numerical control (claim) & cell & change & tolerance & result \\
    \midrule
    smearing MV $\to$ Gaussian (C034) & \cell{aIWO_1p0_s1} & max $|\Delta F| = \num{2.4e-4}$\elevenRyb
      & \num{1e-3}\elevenRyb & pass (a screen) \\
    split species labels (C035) & \cell{aIO_1p0_s1} & $|\Delta E| = \num{1e-8}$\,Ry & \num{1e-5}\,Ry
      & pass \\
    fresh restart, same input (C036) & \cell{aIWO_0p8_s2} & max $|\Delta F| = \num{4.99e-3}$\elevenRyb & --
      & not reproduced; cause open \\
    \bottomrule
  \end{tabularx}\par\endgroup
  \src{PAW-PBE, one geometry each \lab{computed}. Since C036, every rescue-route cell needs a restart gate and a
    qualification SCF ($10^{-5}$\,Ry, $10^{-3}$\elevenRyb); otherwise its relaxation is exploratory.}
\end{frame}

% ---------------------------------------------------------------------------------------------------------------------
\begin{frame}{H-N classifies each W cell by frozen observables}
  \elevendense
  \begin{columns}[T,onlytextwidth]
    \begin{column}{0.50\textwidth}
      {\footnotesize Final SCF, bands $n > N_v$, $\elevenNex = n_{\mathrm{W}}$ excess electrons:\par}
      \vspace{-1.6mm}
      {\footnotesize
      \begin{align*}
        \elevenfext &= \tfrac{1}{\elevenNex}\textstyle\sum f_{n\mathbf{k}}w_{\mathbf{k}}\,
          \mathbf{1}\!\left[\PRn \ge 0.20,\ \varepsilon_{n\mathbf{k}} \ge \elevenEh\right]\\
        \elevenflocW &= \tfrac{1}{\elevenNex}\textstyle\sum f_{n\mathbf{k}}w_{\mathbf{k}}\,
          \mathbf{1}\!\left[\PRn \le 0.05,\ p_{\mathrm{W5d}} \ge 0.30\right]\\
        \elevenDW &= \min\{\varepsilon_{n\mathbf{k}} : p_{\mathrm{W5d}} \ge 0.30\} - \elevenEh
      \end{align*}\par}
      \vspace{-1.2mm}
      {\scriptsize $f, w$: occupation, $k$ weight; $p_{\mathrm{W5d}}$: W-5d weight; $\elevenEh$: host edge, lowest
        state with $\PRn \ge 0.20$ and interior weight (W weight not used).\par}
      \vspace{1.0mm}
      {\footnotesize Participation ratio, L\"owdin atom weights $p_i$, $N$ atoms\srefs{BellDean1970,Thouless1974}:\par}
      \vspace{-1.6mm}
      {\footnotesize\[ \PRn = \frac{1}{N\sum_i p_i^2}, \qquad \elevenNeff = N\,\PRn \]\par}
      \vspace{-1.2mm}
      {\scriptsize Equal weight on $M$ atoms: $\PRn = M/N$. Extended $\PRn \ge 0.20$, localised $\PRn \le 0.05$
        (absolute); $\elevenNeff$ reported for every state.\par}
    \end{column}
    \begin{column}{0.46\textwidth}
      \begingroup\scriptsize
      \renewcommand{\arraystretch}{1.12}
      \begin{tabularx}{\linewidth}{@{}l >{\raggedright\arraybackslash}X@{}}
        \toprule
        class & rule, per W cell \\
        \midrule
        \textsc{donor} & $\elevenfext \ge 0.80$, $\elevenDW > 0.3$\,eV, $\elevenflocW < 0.10$ \\
        \textsc{ecb} & $\elevenfext < 0.50$ or $\elevenflocW \ge 0.50$ \\
        \textsc{partial} & otherwise \\
        \textsc{unresolved} & completeness $< 0.90$, electron count off $> 0.02$\,e, too few bands, qualification
          pending \\
        \bottomrule
      \end{tabularx}\par\endgroup
      \src{\textsc{ecb}: extended-carrier efficacy breakdown.}
      \begin{itemize}
        \item \textsc{ecb} need not mean W stopped donating: a host or surface trap may hold the electron.
          W-centred needs $\ge 30$\,\% trap weight on W and its shell.
        \item Neutral PBE shows where the electron goes in the cell, not a charge-transition
          level\srefs{Freysoldt2014,Lany2008}. PBE favours delocalisation.
      \end{itemize}
    \end{column}
  \end{columns}
  \vspace{0.3mm}
  \src{Register section 1 (A16, A24) and amendment A27 (2026-10-11, 09:09Z), before any slab electronic result. H-N
    is exploratory: read-out rules frozen, forecasts not scored.}
\end{frame}

% ---------------------------------------------------------------------------------------------------------------------
\begin{frame}{Absolute PR thresholds mean different supports per cell}
  \elevendense
  \begin{columns}[T,onlytextwidth]
    \begin{column}{0.50\textwidth}
      \begingroup\footnotesize
      \setlength{\tabcolsep}{4.5pt}
      \renewcommand{\arraystretch}{1.15}
      \begin{tabular}{@{}l r c c c@{}}
        \toprule
        W cell & $N$ & localised & extended & $\PRn$ at \\
               &     & $M \le 0.05N$ & $M \ge 0.20N$ & $M = 7$ \\
        \midrule
        bulk & 126 & 6.3 & 25.2 & 0.056 \\
        1.0\,nm & 138 & 6.9 & 27.6 & 0.051 \\
        0.8, 0.7\,nm & 147 & 7.4 & 29.4 & 0.048 \\
        2.0\,nm, s2 & 264 & 13.2 & 52.8 & 0.027 \\
        2.0\,nm, s1 & 270 & 13.5 & 54.0 & 0.026 \\
        \bottomrule
      \end{tabular}\par\endgroup
      \src{Support size $M$ in atoms implied by the frozen thresholds; $M = 7$: weight spread evenly over W and its
        six first-shell O. From the rule and the atom counts only \lab{derived} (M005).}
    \end{column}
    \begin{column}{0.46\textwidth}
      \begin{itemize}
        \item For a fixed support of $M$ atoms, $\PRn = M/N$ falls as $1/N$: one absolute threshold is a different
          support size in each cell.
        \item An even W + 6\,O state ($M = 7$) is localised at 0.7, 0.8 and 2.0\,nm, but intermediate in bulk and at
          1.0\,nm \lab{derived}.
        \item More concentrated states are localised everywhere: half the weight on W, the rest on six O, gives
          $\elevenNeff = 3.4$.
        \item The rule is unchanged; the table is a reading aid for the released results.
      \end{itemize}
    \end{column}
  \end{columns}
  \takeaway{Localised means at most about 6--7 atoms in bulk and thin cells, about 13 at 2.0\,nm.}
\end{frame}

% ---------------------------------------------------------------------------------------------------------------------
\begin{frame}{By its frozen rules, H-N cannot return a crossover verdict}
  \elevendense
  \begin{columns}[T,onlytextwidth]
    \begin{column}{0.54\textwidth}
      \begingroup\scriptsize
      \setlength{\tabcolsep}{3.4pt}
      \renewcommand{\arraystretch}{1.08}
      \begin{tabularx}{\linewidth}{@{}l c c >{\raggedright\arraybackslash}X@{}}
        \toprule
        control & 2.0\,nm & 0.7\,nm & what it tests \\
        \midrule
        C1 paired substitution & -- & req. & W + O in undoped parent \\
        C2 directional dilution & -- & cut & W halved along one axis \\
        C3 occupation width & req. & req. & 0.001\,Ry reference \\
        C4 Gaussian polish & req. & req. & MV geometry re-relaxed \\
        C5 hybrid functional & -- & cut & HSE06, fixed geometry \\
        C6 spin & req. & req. & nspin = 2, two starts \\
        \bottomrule
      \end{tabularx}\par\endgroup
      \src{As registered (A27) and amended (A37: C2, C5 cut).}
      \vspace{0.4mm}
      \begin{itemize}
        \item \textsc{crossover}: bulk and \emph{both} 2.0\,nm seeds \textsc{donor}, 0.7\,nm \textsc{ecb}, controls
          passing at both endpoints.
        \item \textsc{no crossover}: \textsc{donor} at bulk, 2.0, 1.0, 0.8 and 0.7\,nm, all resolved, controls
          passing at 0.7\,nm.
      \end{itemize}
    \end{column}
    \begin{column}{0.42\textwidth}
      {\footnotesize\textbf{Proposition} (M001). Under the frozen rules and A37, the registered H-N outcome is
        \textsc{unresolved} in every branch.\par}
      \vspace{1.2mm}
      {\footnotesize\textit{Proof.} Both verdicts need both 2.0\,nm seeds \textsc{donor} with C4. Seed 2
        (\cell{aIWO_2p0_s2}) never converged in its only MV segment and continues on the one-attempt Gaussian route (A36): it has no MV
        parent to polish, so C4 cannot be met. No other 2.0\,nm W seed exists; an amendment admitting a native-Gaussian seed
        was declined as result-exposed (A36 erratum). A37 sets the outcome to \textsc{unresolved}; per-cell classes stay
        intact.\hfill$\square$\par}
      \vspace{1.0mm}
      {\scriptsize\color{inkgrey} A37: 2026-10-13, 00:10Z (D-42), before any H-N class was computed; no sealed value
        is used \lab{derived}.\par}
    \end{column}
  \end{columns}
  \takeaway{Neither ``W stops donating below a critical thickness'' nor its opposite can be registered.}
\end{frame}

% ---------------------------------------------------------------------------------------------------------------------
\begin{frame}{Per-cell results stay sealed until D-43's conditions hold}
  \elevendense
  \begin{columns}[T,onlytextwidth]
    \begin{column}{0.47\textwidth}
      {\footnotesize Release decision D-43 (2026-10-14, 10:35Z) is conditional. No read-out runs until:\par}
      \vspace{-0.8mm}
      \begin{enumerate}
        \item the read-out manifest (F10) is complete;
        \item lineage and qualification are audited;
        \item checker F9 returns \texttt{CLAIMS\_OK} for every certificate;
        \item a release clearance is logged.
      \end{enumerate}
      {\footnotesize At commit 3fe6627 no clearance is logged. The outputs are excluded from the public repository
        and were not opened for this talk.\par}
      \vspace{1.6mm}
      \sealed{per-cell H-N classes and controls; H-G, H-H, H-I, H-K, H-M scores}
    \end{column}
    \begin{column}{0.49\textwidth}
      \begingroup\scriptsize
      \renewcommand{\arraystretch}{1.1}
      \begin{tabularx}{\linewidth}{@{}>{\raggedright\arraybackslash}p{2.6cm}>{\raggedright\arraybackslash}X@{}}
        \toprule
        class at 0.7\,nm & will be read as \\
        \midrule
        \textsc{donor}; C3, C4, C6 pass & W electron in extended states of this network; does not show
          donation in the devices \\
        \textsc{ecb}, W-centred (C1) & candidate W-specific loss; without C2, C5 possibly a dilution or functional
          artefact \\
        \textsc{ecb}, other & host or surface trapping, not W chemistry \\
        \textsc{partial}, \textsc{unresolved} & no statement about W \\
        \bottomrule
      \end{tabularx}\par\endgroup
      \src{Reading guide fixed 2026-10-14 without sight of any amorphous electronic quantity.}
      \vspace{0.4mm}
      \begin{itemize}
        \item Not possible: a crossover or its absence (A37), or a thickness trend: one network per thickness is
          not a dependence on $t$.
      \end{itemize}
    \end{column}
  \end{columns}
  \takeaway{A decisive H-N needs two independent 2.0\,nm W seeds in one smearing family, a polish-free endpoint
    rule, and C2, C5.}
\end{frame}
